% ============================================================
%  01_abstract.tex
%  Template spec:
%   - REZUMAT heading: Arial 20 pt Bold Uppercase Center
%   - 3 blank 12pt lines BEFORE heading
%   - 2 blank 12pt lines AFTER heading, then text begins
%   - Text: Arial 12 pt, up to one page
%   - English abstract on page 3
% ============================================================

\thispagestyle{plain}
\selectlanguage{romanian}

\vspace*{3\baselineskip}

\begin{center}
{\fontsize{20}{24}\selectfont\textbf{\MakeUppercase{Rezumat}}}
\end{center}

\vspace{2\baselineskip}

Această lucrare descrie proiectarea și implementarea \airnest{}, un prototip de stație de
monitorizare a calității aerului din interior, bazat pe IoT. Componenta principală a
dispozitivului este un \RPi{}, conectat la cinci senzori de gaz din seria MQ, care includ
senzori pentru monoxid de carbon, hidrogen, GPL/Butan, alcool/Benzen și eCO\textsubscript{2}.
Pe lângă acestea, există un senzor DHT11 pentru măsurarea temperaturii și umidității, un
senzor de sunet KY-037 și un convertor analog-digital MCP3208.

Valorile calculate sunt înregistrate la fiecare 30 de minute. Datele colectate sunt înregistrate
mai întâi într-un fișier CSV local pe Raspberry Pi. Acestea sunt încărcate simultan pe un server
Flask pe AWS EC2 prin intermediul bazei de date PostgreSQL pe AWS RDS. Site-ul web, creat pentru
această stație, oferă mijloace pentru a vizualiza toate citirile prin intermediul graficelor,
a descărca datele brute și a vedea predicțiile făcute de o conductă de predicție prin
învățare automată pentru următoarele 7 zile. Predicțiile sunt făcute folosind date
meteorologice istorice Open-Meteo, măsurători oficiale ale calității aerului din SEE de la
patru stații de monitorizare din Timișoara și citirile live ale senzorilor.

Testele au arătat că măsurătorile efectuate de senzori au rezultate consistente atunci când sunt
efectuate consecutiv și după un anumit timp în care se pot supraîncălzi. S-a stabilit că
întreaga conductă de date end-to-end, procedurile de autentificare, controlul accesului,
filtrarea datelor, cadrul predictiv și modul de întreținere au fost funcționale după mai
multe teste reușite. Sistemul și-a îndeplinit funcția preconizată de a crea o stație de
monitorizare a calității aerului din interior la prețuri accesibile, cu acces permanent la date.

% ── English Abstract: page 3 ────────────────────────────────
\newpage
\thispagestyle{plain}
\selectlanguage{english}

\vspace*{3\baselineskip}

\begin{center}
{\fontsize{20}{24}\selectfont\textbf{\MakeUppercase{Abstract}}}
\end{center}

\vspace{2\baselineskip}

This paper describes the design and implementation of \airnest{}, an IoT indoor air quality
monitoring station prototype. The device's main component is a \RPi{} that is linked up with five
gas sensors of the MQ series, which include sensors for carbon monoxide, hydrogen, LPG/Butane,
alcohol/Benzene and eCO\textsubscript{2}. Besides them there is a DHT11 sensor for measuring
temperature and humidity, a KY-037 sound sensor, and an MCP3208 analog-to-digital converter.

Computed values are recorded every 30 minutes. Collected data is first recorded into a local CSV
file on the Raspberry Pi. It is simultaneously uploaded to a Flask server on AWS EC2 through a
PostgreSQL database on AWS RDS. The website, which was made for this station, provides means to
see all the readings through charts, download the raw data and see the predictions made by a
machine learning prediction pipeline for the next 7 days. Predictions are made using historical
Open-Meteo weather data, official EEA air quality measurements from four Timișoara monitoring
stations and the live sensor readings.

Testing showed that the measurements made by the sensors have consistent readings when made
consecutively and after some time in which they can heat up. It has been established that the
entire end-to-end data pipeline, the authentication procedures, access control, data filtering,
the predictive framework, and the maintenance mode were all functional after multiple successful
tests. Through all of this, the system manages to create an affordable indoor air quality
monitoring station with permanent access to the data.

\newpage
